\section{Residual incidence guide}
\label{sec:residual}

The two constructions named in the limitations paragraph of \citet{whitaker2025} do not supply the density in the denominator of Equation~\eqref{eq:weight}. Equation (19) of \citet{golightly2019} samples a Markov jump process, and the weight (3) is the jump-process likelihood ratio under a Gaussian observation of the state. Theorem 4.2 of \citet{corstanje2023} conditions an inhomogeneous Poisson process on an exact integer end-point, and the ratio (16) contains the transition probability \(h(0, x_0)\). Corollary 4.3.6 and Theorem 4.4.2 of \citet{corstanje2024} are path-space ratios for a guided chemical reaction process conditioned on an exact end constraint. Equations (15) and (25) of \citet{whitaker2017} are Gaussian densities of a diffusion skeleton. We do not put any of those expressions in place of \(q\).

The residual guide keeps the Poisson draw and keeps \(q\) as the product of Poisson probabilities on the same incidence increments. The numerator of Equation~\eqref{eq:weight} is unchanged. The Poisson means change.

At a substep, the linear-noise approximation of cumulative incidence is restarted at the current prevalence. The incidence coordinate \(\eta\) starts at zero, the fundamental matrix \(G\) starts at the identity, and the residual variance \(V\) starts at zero. The contact rate is held at its current value. We integrate equations (4), (8) and (9) of \citet{golightly2023} over the next substep and over the time remaining until the observation. The Jacobian is the matrix printed after their equation (6). The restart is the one they describe after their equation (19). \citet[Section 4.2]{fearnhead2014} state that restart: the ordinary differential equation is re-initialised at the beginning of each interval. The integration uses a classical Runge--Kutta step of length 0.1. That length is a numerical choice, fixed before the comparison in Section~\ref{sec:simulation}. It is not a step in \citet{golightly2023}.

The mean of the binomial observation then uses the linear-noise mean of the remaining infections in place of the frozen product \(h_{\mathrm{inf}}(t_i - \tau)\). The variance remains Equation~\eqref{eq:hstar}'s binomial variance, \(\sigma^{2} = (1 - \rho)\hat{\mu}\). For an observation of infections this is the variance in equation (22) of \citet{golightly2023}.

Let \(R\) denote the residual incidence about \(\eta\), with variance \(V\) after the restart. The covariance between the residual at the next substep and the residual at the end of the remaining interval is \(V_{\Delta\tau}(G_{t_i - \tau} G_{\Delta\tau}^{-1})'\). \citet[Section 3.1.2]{whitaker2017} write the same block as \(P_t\psi_t P_T'\) after their display (24). The observation loads on the infection component by \(\rho\). The conditional mean of the infection residual at the next substep is that infection covariance, times \(\rho(y - \hat{\mu})\), divided by \(\rho^{2} V_{11}(t_i - \tau) + \sigma^{2}\). The removal component is not shifted, as in the scalar form of Equation~\eqref{eq:hstar}.

The Poisson mean for the next infection increment is the ordinary-differential increment of \(\eta\) plus that conditional residual mean. The Poisson mean for the removal increment is the ordinary-differential increment of the removal coordinate alone. A negative mean is not replaced by zero. The weight is zero if the integration stops because a prevalence on the ordinary differential equation would be negative, if the determinant of \(G_{\Delta\tau}\) is smaller than \(10^{-12}\) in absolute value, if either Poisson mean is negative or is not finite, or if the updated counts would be negative. The truncation in Section 3.2 of \citet{golightly2019} is not used.

Equation (6) of \citet{golightly2018} is the constant-hazard correction that \citet{whitaker2025} adapt. It is not this proposal. \citet{fintzi2022} reparameterise a linear-noise approximation for incidence counts. The equations integrated here are those printed by \citet{golightly2023}, which cite that reparameterisation. We do not transcribe an unread display from \citet{fintzi2022} into \(q\).
